% PAGE-BUDGET TARGET: 5 pages
% HARD MAX: 6 pages
\section{Public proof spine and packet reduction}\label{sec:packet-reduction}

The proof burden of \cref{thm:main} is finite at theorem level. More precisely, the result reduces to ten public packets together with one auxiliary weak-window certificate, and every load-bearing step lies on one of two named chains: Route~1 for the mass-gap transport and Lane~A for the sharp-local construction. This section records that reduction explicitly so that the rest of the paper can be read as a controlled packet stack rather than as a synchronized 368-page suite. Appendix~A records the frozen live dependency DAG, and Appendix~C records the master theorem index and canonical numbering crosswalk fixed in Phase~I.

\subsection{Ten-packet reduction and public theorem docket}

\begin{proposition}[Packet reduction of the main theorem]\label{prop:packet-reduction}
\Cref{thm:main} follows from ten public theorem packets together with one auxiliary certification theorem. More precisely:
\begin{enumerate}[label=\textup{(\roman*)}]
    \item Packets~1--2 imply part~\textup{(A)} of \cref{thm:main}.
    \item Packets~4--7 construct the continuum sharp-local state, the full inductive-union sharp-local algebra, and the bounded-base cyclicity needed for part~\textup{(C)}.
    \item Packet~3 together with Packet~8, with the Lane~A seam nodes consumed inside Packet~8, yield part~\textup{(D)}.
    \item Packet~9 packages the downstream Euclidean clustering reformulation together with OS/Wightman existence, field correspondence, and non-triviality.
    \item Packet~10 fixes the exact theorem-level local-net endpoint.
\end{enumerate}
The only additional certification theorem not on this public docket is the weak-window sufficiency certificate of Section~\ref{sec:route1}.
\end{proposition}

\begin{proof}
This is a bookkeeping reduction. Each implication is exactly the package-level theorem statement named in Appendix~A. No additional analytic input enters here. The auxiliary certificate records that every surviving weak-window use factors through the repaired curvature-sector package and that no hidden branch re-enters the public Route~1 chain.
\end{proof}

\subsection{Three first-pass theorem modules}

A first-pass assessment can be bounded by three theorem modules.
\begin{description}[leftmargin=2.2em,style=nextline]
\item[Ultraviolet/public-scope gate.] Packet~1 together with Packet~2. This module decides whether the compact-simple A1 input, public group-scope export, and one-shot entrance theorem are present on the public proof spine.
\item[Route~1 lattice-side burden.] Packet~3 together with the same-scale Wilson return, the Lane~A seam nodes consumed in Packet~8, and the sole live weak-window certificate. This module decides whether a tuned full fixed-lattice lower bound is attached to the physical Wilson theory before the continuum transport begins.
\item[Lane~A and QE3 closure chain.] Packets~4--8. This module decides whether the proof reaches the full sharp-local inductive union, preserves bounded-base cyclicity there, and transports the fixed-lattice lower bound to the continuum sharp-local vacuum sector.
\end{description}

\subsection{Decisive seam nodes}

After packet compression, the reader need only keep the following seam nodes explicitly in view.
\begin{description}[leftmargin=2.2em,style=nextline]
\item[Theorem~F.331(i).] Unique downstream mixed-channel closure node at finite cap.
\item[Promoted bridge into Theorem~5.74.] The theorem-facing positive/unital finite-cap bridge packages exactly what the capwise extension theorem may use from the bounded-family and mixed-channel closure package, and nothing more.
\item[Theorem~5.74(b).] Finite-cap sharp-local OS structure on the capwise algebra; this is the capwise OS input used before any inductive-union passage.
\item[Corollary~5.75.] Bounded-state restriction compatibility node on the QE3 seam; it is \emph{not} the inductive-union theorem.
\item[Corollary~5.76.] Passage from capwise sharp-local reconstruction to the full inductive-union sharp-local algebra.
\item[Theorem~5.77.] First theorem-level bounded-base cyclicity export on the Lane~A spine.
\item[Lemma~F.214.] Final QE3 density/graph-core handoff; closes the cap/flow-time seam with no hidden domain argument.
\item[Theorem~F.216.] Unique transport node carrying the tuned full fixed-lattice lower bound to the continuum sharp-local Hilbert space.
\item[Theorem~F.298.] Sole live weak-window sufficiency certificate on the public Route~1 spine.
\end{description}

\subsection{Route~1 and Lane~A reading order}

The public theorem spine consists of two named load-bearing chains. At the claim-facing level the shortest mass-gap compression is
\[
  \text{Packet 1}\Rightarrow \text{Packet 2}\Rightarrow \text{Packet 3}\Rightarrow \text{Packet 8},
\]
with the proviso that Packet~8 explicitly consumes the Lane~A seam nodes Theorem~5.74(b), Corollary~5.75, Corollary~5.76, Theorem~5.77, Corollary~F.213, and Lemma~F.214, and still requires the weak-window sufficiency certificate of Theorem~F.298. The unique constructive chain is
\[
  \text{Packet 4}\Rightarrow \text{Packet 5}\Rightarrow \text{Packet 6}\Rightarrow \text{Packet 7},
\]
where the mixed-correlator closure node and the promoted finite-cap state bridge occur first, the capwise OS-structure theorem is still Theorem~5.74, the inductive-union passage occurs only afterward at Corollary~5.76, and bounded-base cyclicity appears only at the final theorem of the chain.

For a concentrated first pass the expensive checks are concentrated in three places only: the compact-simple ultraviolet/public-scope gate in Packet~1, the fixed-lattice-to-continuum transport seam across Packets~3 and~8, and the exact-dimension-to-inductive-union closure route across Packets~5--7. Everything else should be treated, on first pass, as certification, endpoint packaging, or explicitly deferred illustration.

From this point onward the reader should track only three front-line modules: the ultraviolet/public-scope gate, the Route~1 fixed-lattice-to-continuum gap chain, and the Lane~A sharp-local closure chain. Everything else is either a supporting proof home, a claim-facing endpoint packet, or deferred illustrative material. The remaining sections state those modules in that same order.
